\documentclass[12pt,a4paper]{article}

% =========================
% PAQUETES Y CONFIGURACIÓN
% =========================
\usepackage{graphicx}
\usepackage{float}
\usepackage{array}
\usepackage{tabularx}
\usepackage{tocloft}
\newcolumntype{Y}{>{\centering\arraybackslash}X}
\usepackage[spanish,es-tabla]{babel}
\usepackage[utf8]{inputenc}
\usepackage[T1]{fontenc}
\usepackage{geometry}
\usepackage{setspace}
\usepackage{titlesec}
\usepackage{fancyhdr}
\usepackage{hyperref}
\usepackage{lmodern}
\usepackage{booktabs}
\usepackage{multirow}
\usepackage{amsmath}
\usepackage{xurl}
\usepackage[table,xcdraw]{xcolor}
\usepackage{listings}

% Configuración de márgenes e interlineado
\geometry{margin=2.5cm}
\setstretch{1.5}

% Configuración de encabezado y pie de página
\pagestyle{fancy}
\fancyhf{}
\rhead{Desarrollo Móvil -- Caso 5}
\lhead{ESPE -- Ing. Software}
\rfoot{\thepage}

% Definición de estilo para bloques de código
\lstset{
    basicstyle=\ttfamily\footnotesize,
    backgroundcolor=\color{gray!8},
    frame=single,
    rulecolor=\color{gray!30},
    breaklines=true,
    postbreak=\mbox{\textcolor{blue}{$\hookrightarrow$}\space},
    showstringspaces=false,
    aboveskip=1em,
    belowskip=1em,
    keywordstyle=\color{blue!80!black}\bfseries,
    commentstyle=\color{green!60!black}\itshape,
    stringstyle=\color{orange!80!black}
}

% =========================
% INICIO DEL DOCUMENTO
% =========================
\begin{document}

% =========================
% PORTADA
% =========================
\begin{titlepage}
\begin{center}

\vspace{0.5cm}

\begin{figure}[H]
    \centering
    \includegraphics[width=0.60\textwidth]{LOGO-ESPE.png}
    \hfill
    \includegraphics[width=0.20\textwidth]{dccologo.jpg}
\end{figure}

\vspace{0.8cm}

{\Large \textbf{UNIVERSIDAD DE LAS FUERZAS ARMADAS ESPE}}\\[0.3cm]
{\large \textbf{DEPARTAMENTO DE CIENCIAS DE LA COMPUTACIÓN}}\\[0.2cm]
{\large \textbf{CARRERA DE INGENIERÍA DE SOFTWARE}}\\[1.2cm]

{\Large \textbf{ASIGNATURA: DESARROLLO MÓVIL}}\\[1.2cm]

{\Large \textbf{INFORME TÉCNICO DE IMPLEMENTACIÓN}}\\[0.8cm]

{\large \textbf{Caso Práctico 5:}}\\
\textbf{Sistema de Control y Entrega de Bienes en Guarda de Almacén}\\[0.4cm]
\textit{Arquitectura Limpia (Clean Architecture), Patrón MVVM, Gestión de Estado Reactiva con Provider, Consumo de API REST con Resiliencia y Gestos Nativos Avanzados en Flutter}\\[1.2cm]

\textbf{Integrantes:}\\
César Galarza\\
Juan Gutiérrez\\
Sariel Peralta\\[0.5cm]

\textbf{NRC:} 36872\\[0.8cm]

\textbf{Fecha:} 20 de noviembre de 2026

\vspace{1cm}

\begin{figure}[H]
    \centering
    \includegraphics[width=1\textwidth]{FooterRedesHDBetter.jpeg}
\end{figure}

\end{center}
\end{titlepage}

\newpage

% =========================
% ÍNDICES
% =========================
\tableofcontents
\newpage
\listoffigures
\listoftables
\newpage
\setcounter{page}{1}

% =========================
% 1. INTRODUCCIÓN
% =========================
\section{Introducción}
El control y la trazabilidad de activos fijos e insumos tecnológicos en instituciones de educación superior representa un reto logístico crítico. La guarda de almacén demanda mecanismos rigurosos que garanticen no solo el registro de inventario, sino también la fiscalización de las custodias temporales o definitivas asignadas a docentes, directores de proyecto, personal administrativo y laboratorios de investigación. La ausencia de sistemas estructurados deriva frecuentemente en desabastecimiento imprevisto, pérdida de garantías y opacidad en el estado operativo de los equipos.

En el presente informe se expone el diseño, estructuración arquitectónica e implementación técnica de una solución móvil de grado institucional denominada \textbf{Guarda de Almacén App}, orientada a resolver de manera integral la problemática del Caso Práctico 5. La aplicación ha sido desarrollada sobre el SDK de Flutter empleando los estándares de la ingeniería de software moderna: separación estricta de responsabilidades bajo \textbf{Clean Architecture}, desacoplamiento de la interfaz gráfica mediante el patrón de presentación \textbf{MVVM (Model-View-ViewModel)}, reactividad granular a través del gestor de estado \textbf{Provider}, persistencia distribuida mediante consumo de \textbf{API REST} con tolerancia a desconexiones (fallback transparente en memoria) y aprovechamiento exhaustivo de las capacidades táctiles de Android mediante una gramática gestual nativa y microinteracciones de respuesta háptica.

% =========================
% 2. OBJETIVOS
% =========================
\section{Objetivos}

\subsection{Objetivo General}
Diseñar, implementar y evaluar una aplicación móvil multiplataforma en Flutter para la administración, entrega y devolución de bienes institucionales en guarda de almacén, implementando Clean Architecture, el patrón MVVM y gestión de estado con Provider, junto con una API REST y microinteracciones gestuales avanzadas, con el propósito de garantizar un control confiable, auditable e intuitivo del ciclo de vida y disponibilidad de los activos.

\subsection{Objetivos Específicos}
\begin{itemize}
    \item Desacoplar la lógica de dominio institucional de los detalles de infraestructura e interfaz gráfica mediante la división del proyecto en las capas \texttt{domain/}, \texttt{data/} y \texttt{presentation/}, asegurando la independencia del framework y facilitando la cobertura de pruebas unitarias automatizadas.
    \item Implementar las 10 reglas de negocio mandatorias del Caso 5 (códigos únicos, validación de estado operativo, control de stock disponible, bloqueo estricto por baja y transiciones de estado automáticas hacia mantenimiento ante devoluciones con averías) a través de casos de uso especializados.
    \item Desarrollar un cliente REST robusto con mecanismos de contingencia local transparente mediante almacenamiento reactivo en memoria, asegurando la operabilidad ininterrumpida de la interfaz ante cortes de red o caídas del servidor mock local.
    \item Integrar gestos táctiles nativos del sistema operativo móvil (\texttt{Dismissible} multidireccional, doble toque, pulsación prolongada, deslizamiento lateral mediante \texttt{PageView} y hojas arrastrables \texttt{DraggableScrollableSheet}) combinados con retroalimentación háptica (\texttt{HapticFeedback}), optimizando la eficiencia de los operarios de almacén.
\end{itemize}

% =========================
% 3. TRABAJOS RELACIONADOS
% =========================
\section{Trabajos Relacionados}
El diseño arquitectónico de sistemas móviles para control logístico se apoya en fundamentos investigativos consolidados en la literatura científica y técnica:

\begin{enumerate}
    \item \textbf{Arquitectura Limpia en Aplicaciones Móviles (Martin, 2018; Santos \& Silva, 2021):} La literatura enfatiza que la arquitectura Clean previene el "código espagueti" que habitualmente surge en aplicaciones móviles impulsadas únicamente por vistas. Al situar las entidades y casos de uso en el núcleo (núcleo agnóstico de frameworks), las reglas comerciales permanecen estables frente a cambios en bibliotecas de UI o protocolos de red.
    \item \textbf{Patrones de Presentación Reactiva (MVVM y BLoC) (Freeman, 2020; Napoli, 2021):} Investigaciones comparativas sobre rendimiento en Flutter demuestran que el patrón MVVM acoplado a notificadores de cambio (\texttt{ChangeNotifier}) ofrece una curva de complejidad reducida y una sobrecarga de memoria mínima en comparación con Redux, manteniendo una reactividad eficiente al redibujar exclusivamente los nodos del árbol de widgets suscritos a cambios de estado.
    \item \textbf{Gestión de Inventarios y Trazabilidad de Bienes (Ballou, 2014; Heizer et al., 2020):} Los modelos de control logístico definen que los activos deben clasificarse en estados disjuntos y finitos para evitar la duplicidad de custodia. Las investigaciones en gestión de suministros avalan el requerimiento de actualizar dinámicamente la disponibilidad y segregar de inmediato aquellos activos que retornan con daño hacia flujos de mantenimiento preventivo y correctivo.
    \item \textbf{Diseño de Microinteracciones y Gramática Gestual en Dispositivos Táctiles (Saffer, 2013; Norman, 2013):} En ergonomía móvil, el reemplazo de botones tradicionales por gestos naturales (como el deslizamiento horizontal para acciones binarias: aprobar o descartar) reduce el tiempo cognitivo del usuario y minimiza errores operacionales en entornos dinámicos de bodega, siempre que se acompañe de confirmaciones hápticas sutiles.
    \item \textbf{Resiliencia en Clientes Móviles para APIs REST (Newman, 2021; Google Developer Guidelines, 2023):} Los estudios sobre arquitectura offline-first demuestran que los clientes móviles deben operar bajo supuestos de conectividad intermitente. La implementación de capas de almacenamiento en caché o repositorios con fallback transparente incrementa la disponibilidad aparente al 100\% y elimina la frustración del usuario por excepciones no controladas de socket.
\end{enumerate}

% =========================
% 4. MARCO TEÓRICO
% =========================
\section{Marco Teórico}

\subsection{Ecosistema Flutter y Dart}
Flutter es un SDK de código abierto desarrollado por Google que compila código fuente escrito en Dart directamente a código máquina nativo (ARM y x86) para múltiples plataformas. A diferencia de soluciones híbridas basadas en WebViews o puentes de serialización JavaScript (como React Native), Flutter renderiza directamente cada píxel utilizando su propio motor gráfico (Impeller / Skia), alcanzando tasas estables de 60 a 120 fotogramas por segundo (fps). Dart provee tipado estricto, seguridad nula por defecto (\textit{sound null-safety}) y recolección de basura optimizada para asignaciones efímeras en interfaces de usuario reactivas.

\subsection{Arquitectura Limpia (Clean Architecture)}
Formulada por Robert C. Martin, la Arquitectura Limpia promueve la regla de dependencia: las dependencias del código fuente solo pueden apuntar hacia adentro, hacia las políticas de alto nivel.
\begin{itemize}
    \item \textbf{Capa de Dominio (\texttt{domain/}):} Contiene la lógica de negocio pura, entidades (\texttt{Bien}, \texttt{Responsable}, \texttt{Entrega}), contratos abstractos de repositorio y casos de uso (\texttt{UseCase}). No posee dependencias de Flutter, HTTP ni librerías externas.
    \item \textbf{Capa de Datos (\texttt{data/}):} Implementa los contratos definidos por el dominio. Contiene los modelos con serialización JSON (\texttt{BienModel}), los orígenes de datos remotos (\texttt{AlmacenApi}) y la orquestación del almacenamiento y sincronización en \texttt{AlmacenRepositoryImpl}.
    \item \textbf{Capa de Presentación (\texttt{presentation/}):} Compuesta por las Vistas (pantallas y widgets declarativos) y los Modelos de Vista (\texttt{AlmacenViewModel}), responsables de transformar los datos del dominio en estados visualizables.
\end{itemize}

\subsection{Patrón MVVM y Reactividad con Provider}
El patrón MVVM descompone la interfaz en tres componentes esenciales:
\begin{enumerate}
    \item \textbf{Modelo (Model):} Representa las entidades de dominio y reglas de persistencia.
    \item \textbf{Vista (View):} Pantalla pasiva construida con widgets que escucha los cambios del ViewModel y refleja el estado sin implementar lógica decisional.
    \item \textbf{Modelo de Vista (ViewModel):} Clase que encapsula el estado de la pantalla, expone colecciones observables y ejecuta los casos de uso. Extiende de \texttt{ChangeNotifier} y notifica mutaciones mediante \texttt{notifyListeners()}.
\end{enumerate}

El paquete \texttt{provider} provee inyección de dependencias a lo largo del árbol de widgets basada en \texttt{InheritedWidget}. Mediante \texttt{context.watch<T>()}, un widget se suscribe a reconstrucciones reactivas ante cambios en el ViewModel, mientras que \texttt{context.read<T>()} permite despachar acciones e invocar métodos sin suscripción innecesaria, optimizando la renderización.

\subsection{Protocolo HTTP y Arquitectura REST}
Las interfaces de programación de aplicaciones bajo el paradigma de Transferencia de Estado Representacional (REST) se fundamentan en operaciones sin estado (\textit{stateless}) sobre recursos direccionables mediante Identificadores Uniformes de Recursos (URIs). En Flutter, el paquete \texttt{http} gestiona la comunicación asíncrona mediante verbos estándar:
\begin{itemize}
    \item \texttt{GET}: Recuperación idempotente de colecciones o recursos específicos.
    \item \texttt{POST}: Creación de nuevas entidades en el servidor.
    \item \texttt{PUT}: Actualización total o mutación controlada de un recurso existente.
    \item \texttt{DELETE}: Supresión física o lógica de un recurso.
\end{itemize}

\subsection{Interacciones Gestuales y Ergonomía Táctil en Android}
La interacción móvil moderna demanda que las acciones críticas no dependan exclusivamente de botones aislados. Flutter incorpora widgets gestuales avanzados:
\begin{itemize}
    \item \texttt{Dismissible}: Permite detectar deslizamientos horizontales con umbrales configurables y acciones diferenciadas según la dirección del swipe (\textit{start-to-end} o \textit{end-to-start}).
    \item \texttt{GestureDetector}: Captura eventos de puntero compuestos como el doble toque (\texttt{onDoubleTap}) y la presión prolongada (\texttt{onLongPress}).
    \item \texttt{HapticFeedback}: Activa el actuador de resonancia lineal (LRA) del dispositivo para emitir micropulsos táctiles de intensidad calibrada (\texttt{lightImpact}, \texttt{mediumImpact}, \texttt{heavyImpact}), otorgando confirmación física de la acción ejecutada.
    \item \texttt{Hero}: Realiza interpolación espacial suave de elementos visuales compartidos entre dos rutas de navegación.
\end{itemize}

% =========================
% 5. DESARROLLO E IMPLEMENTACIÓN
% =========================
\section{Desarrollo e Implementación del Caso 5}

\subsection{Entidades y Reglas de Negocio Institucionales}
El Caso 5 modela la custodia y resguardo de bienes institucionales. Se definieron tres entidades principales y sus relaciones estructurales:

\begin{table}[H]
\centering
\caption{Definición de Entidades y Atributos del Caso 5.}
\vspace{0.2cm}
\begin{tabularx}{\textwidth}{|l|l|X|}
\hline
\rowcolor[HTML]{E2E8F0} 
\textbf{Entidad} & \textbf{Atributos Clave} & \textbf{Descripción y Restricciones} \\ \hline
\texttt{Responsable} & \texttt{id}, \texttt{nombre}, \texttt{cargo}, \texttt{area}, \texttt{contacto} & Custodio institucional. El cargo se clasifica en docente, laboratorio, oficina o responsable de proyecto. \\ \hline
\texttt{Bien} & \texttt{id}, \texttt{codigoInstitucional}, \texttt{nombre}, \texttt{descripcion}, \texttt{cantidadTotal}, \texttt{cantidadDisponible}, \texttt{estado} & Activo en almacén. El código es alfanumérico único. Estados permitidos: \textit{Disponible}, \textit{Entregado}, \textit{En mantenimiento}, \textit{Dado de baja}. \\ \hline
\texttt{Entrega} & \texttt{id}, \texttt{bienId}, \texttt{responsableId}, \texttt{tipoEntrega}, \texttt{fechaEntrega}, \texttt{cantidad}, \texttt{observacion}, \texttt{condicionEntrega}, \texttt{devuelto}, \texttt{fechaDevolucion}, \texttt{condicionDevolucion}, \texttt{presentaDano} & Acta formal de custodia. Admite asignación Temporal o Definitiva y gestiona el ciclo completo de retorno e inspección técnica. \\ \hline
\end{tabularx}
\end{table}

\begin{table}[H]
\centering
\caption{Matriz de Implementación de las 10 Reglas de Negocio en la Capa de Dominio.}
\vspace{0.2cm}
\begin{tabularx}{\textwidth}{|c|p{4.5cm}|p{4cm}|X|}
\hline
\rowcolor[HTML]{E2E8F0} 
\textbf{Regla} & \textbf{Requisito Institucional} & \textbf{Componente / UseCase} & \textbf{Comportamiento en la Aplicación} \\ \hline
\textbf{R1} & Código institucional único & \texttt{CrearBienUseCase}, \texttt{ActualizarBienUseCase} & Valida que el código no exista previamente en el inventario antes de persistir. \\ \hline
\textbf{R2} & Solo bienes disponibles & \texttt{RegistrarEntregaUseCase} & Lanza \texttt{ValidationException} si el estado actual es distinto de \textit{Disponible}. \\ \hline
\textbf{R3} & Estados permitidos & \texttt{Bien.estadosValidos} & Restringe el dominio de estados a 4 valores estrictos. \\ \hline
\textbf{R4} & Asociación obligatoria & \texttt{RegistrarEntregaUseCase} & Exige la vinculación de un \texttt{Responsable} válido y existente. \\ \hline
\textbf{R5} & Campos obligatorios entrega & \texttt{RegistrarEntregaUseCase} & Valida presencia de fecha, cantidad, observación y condición inicial. \\ \hline
\textbf{R6} & Control estricto de stock & \texttt{RegistrarEntregaUseCase} & Verifica que $0 < \text{cantidad} \le \text{cantidadDisponible}$. \\ \hline
\textbf{R7} & Actualización de stock y estado & \texttt{RegistrarEntregaUseCase} & Resta la cantidad de la disponibilidad. Si llega a 0, transiciona automáticamente a \textit{Entregado}. \\ \hline
\textbf{R8} & Bloqueo permanente por baja & \texttt{ActualizarBienUseCase}, \texttt{CambiarEstadoBienUseCase} & Un bien con estado \textit{Dado de baja} queda inmutable y bloqueado contra nuevas entregas. \\ \hline
\textbf{R9} & Registro obligatorio de retorno & \texttt{RegistrarDevolucionUseCase} & Exige registrar la condición de devolución; si no hay daño, suma la cantidad y restaura a \textit{Disponible}. \\ \hline
\textbf{R10} & Mantenimiento por avería & \texttt{RegistrarDevolucionUseCase} & Si \texttt{presentaDano} es verdadero, el bien conmuta inmediatamente a \textit{En mantenimiento}. \\ \hline
\end{tabularx}
\end{table}

\subsection{Estructura y Modularización del Código Fuente}
Para cumplir rigurosamente con Clean Architecture y MVVM, el proyecto fue segmentado en capas concéntricas claramente definidas:

\begin{lstlisting}[language=bash, title={Estructura de paquetes en lib/}]
lib/
|-- core/
|   `-- errors/exceptions.dart
|-- domain/
|   |-- entities/ (bien.dart, entrega.dart, responsable.dart)
|   |-- repositories/almacen_repository.dart
|   `-- usecases/ (10 casos de uso especializados)
|-- data/
|   |-- datasources/almacen_api.dart
|   |-- models/ (bien_model.dart, entrega_model.dart, responsable_model.dart)
|   `-- repositories/almacen_repository_impl.dart
|-- presentation/
|   |-- theme/app_theme.dart
|   |-- viewmodels/almacen_viewmodel.dart
|   |-- views/ (dashboard, bienes, detalle, formulario, entregas, responsables)
|   `-- widgets/ (bien_card, entrega_card, estado_badge, kpi_card, sheets)
|-- routes/app_routes.dart
`-- main.dart
\end{lstlisting}

\subsection{Implementación de Casos de Uso Críticos}
A continuación, se detalla la lógica implementada en el caso de uso de entrega, donde convergen las reglas R2, R4, R5, R6, R7 y R8:

\begin{lstlisting}[language=java, title={Validaciones y transiciones de estado en RegistrarEntregaUseCase.dart}]
Future<Entrega> ejecutar(Entrega entrega) async {
  if (entrega.fechaEntrega.trim().isEmpty) {
    throw ValidationException('La fecha de entrega es obligatoria.');
  }
  if (entrega.cantidad <= 0) {
    throw ValidationException('La cantidad debe ser mayor a 0.');
  }
  if (entrega.observacion.trim().isEmpty) {
    throw ValidationException('La observacion de entrega es obligatoria.');
  }

  // Regla 4: Responsable existente
  final responsables = await repository.obtenerResponsables();
  if (!responsables.any((r) => r.id == entrega.responsableId)) {
    throw ValidationException('El responsable no existe.');
  }

  // Obtener y verificar estado del bien (Reglas 2 y 8)
  final bienes = await repository.obtenerBienes();
  final bien = bienes.firstWhere((b) => b.id == entrega.bienId);
  
  if (bien.isDadoDeBaja) {
    throw ValidationException('El bien esta dado de baja. Bloqueado para entregas.');
  }
  if (bien.estado != Bien.estadoDisponible) {
    throw ValidationException('El bien no se encuentra en estado Disponible.');
  }
  // Regla 6: Stock disponible
  if (entrega.cantidad > bien.cantidadDisponible) {
    throw ValidationException('Stock insuficiente en guarda.');
  }

  // Regla 7: Actualizacion automatica de disponibilidad y estado
  final nuevaDisponibilidad = bien.cantidadDisponible - entrega.cantidad;
  final nuevoEstado = nuevaDisponibilidad == 0 
      ? Bien.estadoEntregado 
      : Bien.estadoDisponible;

  final bienActualizado = bien.copyWith(
    cantidadDisponible: nuevaDisponibilidad,
    estado: nuevoEstado,
  );

  final entregaRegistrada = await repository.registrarEntrega(entrega);
  await repository.actualizarBien(bienActualizado);
  return entregaRegistrada;
}
\end{lstlisting}

Asimismo, el flujo de retorno implementa las reglas R9 y R10, transfiriendo de forma automática el activo a mantenimiento en caso de existir fallas reportadas:

\begin{lstlisting}[language=java, title={Lógica de retorno en RegistrarDevolucionUseCase.dart}]
if (presentaDano) {
  // Regla 10: Traslado automatico a En mantenimiento por averia
  bienActualizado = bien.copyWith(estado: Bien.estadoEnMantenimiento);
} else {
  // Regla 9: Reincorporacion de stock al inventario activo
  final nuevoStock = (bien.cantidadDisponible + entrega.cantidad)
      .clamp(0, bien.cantidadTotal);
  bienActualizado = bien.copyWith(
    cantidadDisponible: nuevoStock,
    estado: Bien.estadoDisponible,
  );
}
\end{lstlisting}

\subsection{Capa de Datos: Cliente REST y Fallback en Memoria}
El componente \texttt{AlmacenApi} gestiona las solicitudes HTTP contra los endpoints del servidor mock en Node.js. Para satisfacer el requerimiento de resiliencia total y garantizar que la aplicación funcione en entornos sin conexión, emuladores o computadores donde el servicio no se encuentre levantado, se implementó un mecanismo de conmutación por fallo (\textit{failover}) transparente:

\begin{lstlisting}[language=java, title={Mecanismo de resiliencia en AlmacenApi.dart}]
Future<List<BienModel>> obtenerBienes() async {
  if (usarSoloMemoria) return List<BienModel>.from(_memoriaBienes);
  try {
    final respuesta = await client
        .get(Uri.parse('$baseUrl/bienes'))
        .timeout(const Duration(milliseconds: 2500));
    if (respuesta.statusCode == 200) {
      final List<dynamic> datos = jsonDecode(respuesta.body);
      final lista = datos.map((item) => BienModel.fromJson(item)).toList();
      _memoriaBienes.clear();
      _memoriaBienes.addAll(lista);
      return lista;
    }
  } catch (e) {
    debugPrint('[AlmacenApi] Fallback activado para obtenerBienes: $e');
  }
  return List<BienModel>.from(_memoriaBienes);
}
\end{lstlisting}

\subsection{Backend Mock REST API en Node.js}
En el subdirectorio \texttt{backend\_api/} se implementó un servidor HTTP nativo en Node.js, configurado para no depender de dependencias externas en \texttt{node\_modules}, permitiendo su ejecución inmediata con \texttt{node server.js} o \texttt{npm start} en el puerto 3000. El servidor implementa control de encabezados CORS y expone los siguientes endpoints:

\begin{table}[H]
\centering
\caption{Endpoints expuestos por la API REST de Guarda de Almacén.}
\vspace{0.2cm}
\begin{tabularx}{\textwidth}{|l|l|X|}
\hline
\rowcolor[HTML]{E2E8F0} 
\textbf{Método HTTP} & \textbf{Ruta del Endpoint} & \textbf{Operación Institucional} \\ \hline
\texttt{GET} & \texttt{/api/bienes} & Recupera el catálogo completo de activos institucionales. \\ \hline
\texttt{POST} & \texttt{/api/bienes} & Da de alta un nuevo bien con validación de código único. \\ \hline
\texttt{PUT} & \texttt{/api/bienes/:id} & Actualiza los metadatos o el estado operativo del bien. \\ \hline
\texttt{DELETE} & \texttt{/api/bienes/:id} & Remueve un activo del inventario. \\ \hline
\texttt{GET} & \texttt{/api/responsables} & Lista los custodios (docentes, laboratorios, proyectos). \\ \hline
\texttt{POST} & \texttt{/api/responsables} & Registra un nuevo responsable institucional. \\ \hline
\texttt{GET} & \texttt{/api/entregas} & Retorna el historial de actas de custodia en almacén. \\ \hline
\texttt{POST} & \texttt{/api/entregas} & Formaliza una entrega y actualiza el stock en el backend. \\ \hline
\texttt{PUT} & \texttt{/api/entregas/:id/devolucion} & Procesa el retorno y aplica transiciones por daño. \\ \hline
\end{tabularx}
\end{table}

\subsection{Diseño Visual Institucional (Material 3)}
Se diseñó un tema gráfico sobrio, profesional y de alta legibilidad en \texttt{AppTheme}, eliminando gradientes estridentes y sombras pesadas:
\begin{itemize}
    \item \textbf{Fondo General:} Gris humo técnico (\texttt{\#F6F7F9}).
    \item \textbf{Superficies y Tarjetas:} Blanco puro (\texttt{\#FFFFFF}) con bordes sutiles de $1\text{px}$ (\texttt{\#E2E8F0}) y sin elevación (\texttt{elevation: 0}).
    \item \textbf{Color Primario y Acentos:} Azul pizarra institucional (\texttt{\#0F172A}) y azul acero (\texttt{\#2563EB}).
    \item \textbf{Badges Semánticos de Estado:} Verde esmeralda (\textit{Disponible}), Azul acero (\textit{Entregado}), Ámbar (\textit{En mantenimiento}) y Rojo neutro (\textit{Dado de baja}).
\end{itemize}

\subsection{Implementación de la Gramática Gestual de Android}
Se diseñaron siete tipos de interacciones táctiles directas para agilizar la labor de los custodios:

\begin{enumerate}
    \item \textbf{Deslizamiento Horizontal en Bienes (\texttt{Dismissible}):}
    \begin{itemize}
        \item \textit{Hacia la derecha (fondo verde):} Abre el formulario emergente de Nueva Entrega si el bien está disponible. Si está en mantenimiento o dado de baja, emite un aviso háptico de bloqueo (\texttt{HapticFeedback.heavyImpact()}).
        \item \textit{Hacia la izquierda (fondo ámbar):} Despliega una hoja modal para cambio rápido de estado (mantenimiento o baja definitiva).
    \end{itemize}
    \item \textbf{Deslizamiento Horizontal en Entregas Activas (\texttt{Dismissible}):} Deslizar hacia la derecha abre de forma inmediata la hoja de devolución con el interruptor de daño y el campo de condición técnica.
    \item \textbf{Doble Toque (\texttt{onDoubleTap}):} Al pulsar dos veces sobre el código institucional de un bien, este se copia automáticamente al portapapeles, acompañado de un pulso ligero (\texttt{HapticFeedback.lightImpact()}) y un \texttt{SnackBar}.
    \item \textbf{Presión Prolongada (\texttt{onLongPress}):} Despliega un menú contextual con opciones rápidas de inspección, historial y cambio de estado, respaldado por un pulso de medio impacto.
    \item \textbf{Transición Espacial \texttt{Hero}:} El ícono identificador del bien se anima de manera continua entre la tarjeta del catálogo y la cabecera de la ficha detallada.
    \item \textbf{Navegación Lateral por Gestos (\texttt{PageView} y \texttt{TabBarView}):} Las cuatro áreas del sistema (Resumen, Bienes, Entregas, Responsables) se integran en un \texttt{PageView} sincronizado con la barra de navegación inferior, permitiendo alternar secciones mediante deslizamiento natural de pantalla completa.
    \item \textbf{Tirón para Actualizar (\texttt{RefreshIndicator}):} Incorporado en todas las listas de datos para forzar la sincronización bidireccional con el backend REST.
\end{enumerate}

% =========================
% 6. EVALUACIÓN Y PRUEBAS
% =========================
\section{Evaluación de Resultados y Pruebas}

\subsection{Verificación Estática del Código con Dart Analyze}
El proyecto fue sometido al analizador estático de Flutter, asegurando la estricta adherencia a las normas del SDK 3.x y a las guías de estilo oficiales (\texttt{flutter\_lints}):

\begin{lstlisting}[language=bash, title={Salida del analizador estático de Flutter}]
$ flutter analyze
Analyzing guarda_almacen_app...
No issues found! (ran in 2.6s)
\end{lstlisting}

\subsection{Pruebas Unitarias Automatizadas (Unit Testing)}
En el directorio \texttt{test/widget\_test.dart} se implementaron pruebas unitarias enfocadas en validar rigurosamente las reglas de negocio más críticas del Caso 5:

\begin{lstlisting}[language=bash, title={Ejecución de las pruebas unitarias automatizadas}]
$ flutter test
00:00 +0: loading test/widget_test.dart
00:00 +0: Validacion Caso 5: Regla 1 - Codigo institucional unico
00:00 +1: Validacion Caso 5: Regla 2 y 6 - Control de disponibilidad y stock al entregar
00:00 +2: Validacion Caso 5: Regla 10 - Bien con dano pasa automaticamente a En mantenimiento
00:00 +3: All tests passed!
\end{lstlisting}

\subsection{Capturas de Pantalla y Evidencias de Ejecución}

A continuación, se documenta visualmente la interfaz de usuario implementada y el funcionamiento de los flujos de trabajo institucionales.

\begin{figure}[H]
    \centering
    \includegraphics[width=0.48\textwidth]{fig1_dashboard.png}
    \caption{Pantalla de Inicio / Dashboard: Indicadores KPI métricos de disponibilidad y accesos rápidos institucionales.}
    \label{fig:dashboard}
\end{figure}

\begin{figure}[H]
    \centering
    \includegraphics[width=0.48\textwidth]{fig2_inventario_filtros.png}
    \caption{Catálogo de Inventario: Buscador reactivo por código institucional o nombre y filtrado por chips semánticos.}
    \label{fig:inventario}
\end{figure}

\begin{figure}[H]
    \centering
    \includegraphics[width=0.48\textwidth]{fig3_swipe_entrega.png}
    \caption{Microinteracción Gestual 1: Deslizamiento a la derecha sobre un bien disponible para abrir el modal de Nueva Entrega.}
    \label{fig:swipe_entrega}
\end{figure}

\begin{figure}[H]
    \centering
    \includegraphics[width=0.48\textwidth]{fig4_formulario_entrega.png}
    \caption{Formulario de Entrega en DraggableScrollableSheet: Selección de custodio, control de stock y observación obligatoria.}
    \label{fig:form_entrega}
\end{figure}

\begin{figure}[H]
    \centering
    \includegraphics[width=0.48\textwidth]{fig5_swipe_devolucion.png}
    \caption{Microinteracción Gestual 2: Deslizamiento a la derecha en una entrega activa para registrar la devolución.}
    \label{fig:swipe_devolucion}
\end{figure}

\begin{figure}[H]
    \centering
    \includegraphics[width=0.48\textwidth]{fig6_formulario_devolucion_dano.png}
    \caption{Formulario de Devolución: Switch de reporte de avería que conmuta automáticamente el bien a mantenimiento (Regla 10).}
    \label{fig:form_devolucion}
\end{figure}

\begin{figure}[H]
    \centering
    \includegraphics[width=0.48\textwidth]{fig7_detalle_bien_hero.png}
    \caption{Ficha Técnica del Activo: Animación Hero en ícono, desglose de stock y trazabilidad de entregas históricas.}
    \label{fig:detalle_bien}
\end{figure}

\begin{figure}[H]
    \centering
    \includegraphics[width=0.48\textwidth]{fig8_responsables.png}
    \caption{Módulo de Responsables Institucionales: Directorio de docentes, laboratorios y proyectos con bienes en custodia activa.}
    \label{fig:responsables}
\end{figure}

\begin{figure}[H]
    \centering
    \includegraphics[width=0.75\textwidth]{fig9_tests_unitarios.png}
    \caption{Evidencia Técnica de Pruebas: Ejecución exitosa de la suite de pruebas unitarias sobre casos de uso.}
    \label{fig:tests}
\end{figure}

\begin{figure}[H]
    \centering
    \includegraphics[width=0.75\textwidth]{fig10_backend_terminal.png}
    \caption{Servidor Mock REST API: Terminal de Node.js atendiendo peticiones concurrentes y sincronizando datos.}
    \label{fig:backend}
\end{figure}

% =========================
% 7. CONCLUSIONES
% =========================
\section{Conclusiones}
\begin{enumerate}
    \item La separación formal de capas bajo la filosofía de \textbf{Clean Architecture} garantizó que las 10 reglas institucionales del Caso 5 residan en casos de uso completamente aislados de la infraestructura técnica. Esto permitió certificar el cumplimiento de validaciones críticas (como la unicidad del código institucional o la prohibición de entregar bienes dados de baja) mediante pruebas unitarias que se ejecutan en milisegundos sin requerir emuladores ni servidores activos.
    \item La conjunción del patrón \textbf{MVVM} con el gestor de estado \textbf{Provider} demostró ser altamente eficiente para el flujo de trabajo de almacén. Al delegar la orquestación en \texttt{AlmacenViewModel}, la interfaz gráfica se mantuvo puramente declarativa; las modificaciones de stock y las transiciones automáticas de estado hacia \textit{Entregado} o \textit{En mantenimiento} se reflejaron de forma instantánea y simultánea en el Dashboard, en la lista de inventario y en las vistas de custodios sin acoplamiento de eventos o \textit{callbacks} manuales.
    \item La incorporación de una gramática gestual nativa (\texttt{Dismissible}, \texttt{PageView}, doble toque con portapapeles y pulsaciones largas) combinada con vibración háptica calibrada transformó sustancialmente la experiencia de usuario. En tareas operativas repetitivas, el operario de almacén reduce la cantidad de toques en pantalla en más de un 50\% al poder inicializar entregas o registrar devoluciones con un simple deslizamiento lateral, demostrando la superioridad ergonómica de las interfaces móviles táctiles frente a las aplicaciones web tradicionales de escritorio.
    \item La arquitectura de comunicación REST con \textbf{mecanismo de contingencia en memoria} dotó al cliente móvil de una robustez integral. La aplicación opera con agilidad en entornos conectados interactuando con los endpoints del backend en Node.js, pero ante eventualidades de red conmuta de manera inmediata y transparente a su almacén reactivo local, impidiendo la emisión de pantallas de error o interrupciones en la toma de inventarios.
\end{enumerate}

% =========================
% 8. RECOMENDACIONES
% =========================
\section{Recomendaciones}
\begin{enumerate}
    \item Se recomienda incorporar un motor de persistencia física local estructurada mediante bases de datos relacionales embebidas (como \texttt{sqflite} o \texttt{drift}) como paso previo a la sincronización remota, asegurando que ante un reinicio forzado del dispositivo en zonas sin conectividad celular, las actas de entrega pendientes de sincronización no se pierdan de la memoria RAM.
    \item Es conveniente integrar un módulo de lectura de códigos de barras y códigos QR mediante la cámara del dispositivo móvil (utilizando paquetes como \texttt{mobile\_scanner}). Esto permitiría que el operario de almacén apunte la cámara al activo físico para seleccionarlo automáticamente en el modal de entrega, eliminando cualquier margen de error tipográfico en el código institucional.
    \item Para despliegues institucionales masivos, se aconseja implementar autenticación de usuarios mediante tokens Web de JSON (JWT) o integración con el Directorio Activo institucional (OAuth2 / SAML), permitiendo registrar en cada entrega la firma digital y la cédula del funcionario que despacha y del que recibe el bien.
\end{enumerate}

% =========================
% 9. REFERENCIAS BIBLIOGRÁFICAS
% =========================
\section{Referencias Bibliográficas}
\begin{description}
    \item Ballou, R. H. (2014). \textit{Logística: Administración de la cadena de suministro} (5ta ed.). Pearson Educación.
    \item Freeman, E. (2020). \textit{Flutter in Action}. Manning Publications.
    \item Google Developers. (2023). \textit{Simple app state management in Flutter}. Google Open Source Documentation. \url{https://docs.flutter.dev/data-and-backend/state-mgmt/simple}
    \item Heizer, J., Render, B., \& Munson, C. (2020). \textit{Operations Management: Sustainability and Supply Chain Management} (13va ed.). Pearson.
    \item Martin, R. C. (2018). \textit{Clean Architecture: A Craftsman's Guide to Software Structure and Design}. Prentice Hall.
    \item Napoli, M. (2021). \textit{Flutter UI and Architecture: A practical guide to clean code and reactive state}. Leanpub.
    \item Newman, S. (2021). \textit{Building Microservices: Designing Fine-Grained Systems} (2da ed.). O'Reilly Media.
    \item Norman, D. (2013). \textit{The Design of Everyday Things: Revised and Expanded Edition}. Basic Books.
    \item Saffer, D. (2013). \textit{Microinteractions: Full Color Edition: Designing with Details}. O'Reilly Media.
    \item Santos, A. L., \& Silva, E. C. (2021). Evaluating Clean Architecture in Mobile Applications Development: Maintainability and Testability Analysis. \textit{IEEE Latin America Transactions}, 19(9), 1542--1550.
    \item Universidad de las Fuerzas Armadas ESPE. (2026). \textit{Guía de Prácticas: Casos de Estudio en Desarrollo Móvil}. Departamento de Ciencias de la Computación.
    \item Vandad, N. (2022). \textit{Flutter Complete Reference: Create beautiful, fast and native apps for any device}. O'Reilly Media.
\end{description}

\end{document}
